\dictchar{अ}

\bhojentry{अ}{\K{\ba}}{A}{स्वर / उपसर्ग | तद्भव}{\starsThree}%
{हिंदी तथा संस्कृत का पहला स्वर वर्ण; अकार; अभाव या निषेध का उपसर्ग (बिना, ना)।}%
{अ (Prefix: बिना / ना) --- जैसे: अकाल, अनभिन।}%
{First vowel of Indic alphabet; negative prefix (un-, non-, in-).}
\bhojexample{अकाल के समय सब लोग परेशान रहे।}{Everyone was distressed during the famine.}

\bhojentry{अकचब}{\K{\ba\bka\bca\bba}}{Akacab}{हि. क्रि. | [ठेठ]}{\starsTwo}%
{असमंजस या घबराहट में पड़ना।}%
{अकचब (Akacab) / घबरावल}%
{To get confused or perplexed.}

\bhojentry{अकड़ब}{\K{\ba\bka\bdda\bnukta\bba}}{AkaRab}{हि. क्रि. | [ठेठ]}{\starsTwo}%
{ऐंठना, घमंड दिखाना।}%
{अकड़ब / ऐंठल}%
{To act arrogant, stiffen up.}

\bhojsubentry{अकड़}{\K{\ba\bka\bdda\bnukta}}{AkaR}{हि. स्त्री. | [ठेठ]}{\starsTwo}%
{घमंड, ऐंठ, कठोरता।}%
{अकड़ / ऐंठ}%
{Arrogance, stiffness, pride.}

\bhojentry{अकासी}{\K{\ba\bka\bmaa\bsa\bmii}}{Akaasi}{हि. स्त्री. | [पूर्वी/कृषि]}{\starsOne}%
{फल तोड़ने की बांस की लगी कांटा।}%
{अकासी (Akaasi) / लग्गी}%
{Pole hook for harvesting fruits.}

\bhojentry{अकाल}{\K{\ba\bka\bmaa\bla}}{Akaal}{हि. पुं. | तद्भव}{\starsThree}%
{सूखा, दुर्भिक्ष, अन्न का संकट।}%
{अकाल (Akaal) / अकाल-दुर्भिक्ष}%
{Famine, drought, scarcity.}
\bhojexample{अकाल पड़ला पर अनाज महँग हो गइल।}{Grains became expensive when famine struck.}

\bhojentry{अकास}{\K{\ba\bka\bmaa\bsa}}{Akaas}{हि. पुं. | तद्भव}{\starsThree}%
{गगन, अंबर, आसमान।}%
{अकास (Akaas) / आसमान}%
{Sky, firmament.}

\bhojentry{अकेला}{\K{\ba\bka\bme\bla\bmaa}}{Akela}{हि. वि. | तद्भव}{\starsThree}%
{एकमात्र, संगहीन, एकाकी।}%
{अकेला (Akela) / एकेगो}%
{Alone, solitary, single.}

\bhojentry{अकोरी}{\K{\ba\bka\bmo\bra\bmii}}{Akori}{हि. स्त्री. | [दक्षिणी]}{\starsOne}%
{रिश्वत, उकसाने की भेंट।}%
{अकोरी (Akori) / घूस}%
{Bribe, token inducement.}

\bhojentry{अक्किल}{\K{\ba\bka\bvir\bka\bmi\bla}}{Akkil}{हि. स्त्री. | अरबी आगत}{\starsThree}%
{बुद्धि, समझ, समझदारी।}%
{अक्किल (Akkil) / समझ}%
{Intelligence, wisdom, sense.}
\bhojexample{लइका में बहुत अक्किल बा।}{The child has great intelligence.}

\bhojentry{अक्षत}{\K{\ba\bka\bvir\bsha\bta}}{सं. पुं. | तत्सम}{\starsTwo}%
{पूजा का अटूट कच्चा चावल।}%
{अक्षत (Akshat) / अक्षत-चावल}%
{Unbroken sacred rice grains.}

\bhojentry{अक्षर}{\K{\ba\bka\bvir\bsha\bra}}{Akshar}{सं. पुं. | तत्सम}{\starsThree}%
{वर्ण, हरफ, ध्वनि प्रतीक।}%
{अक्षर (Akshar) / आखर}%
{Letter, character, syllable.}

\bhojentry{अखरब}{\K{\ba\bkha\bra\bba}}{Akharab}{हि. क्रि. | [ठेठ]}{\starsTwo}%
{बुरा लगना, खटकना।}%
{अखरब (Akharab) / खटकल}%
{To displease, offend, feel bad.}

\bhojentry{अखड़ा}{\K{\ba\bkha\bdda\bnukta\bmaa}}{AkhaDa}{हि. पुं. | तद्भव}{\starsThree}%
{कुश्ती लड़ने का स्थान, व्यायामशाला।}%
{अखड़ा (AkhaDa) / अखाड़ा}%
{Wrestling arena, gymnasium.}

\bhojentry{अखबार}{\K{\ba\bkha\bba\bmaa\bra}}{Akhbaar}{हि. पुं. | अरबी आगत}{\starsThree}%
{समाचार पत्र।}%
{अखबार (Akhbaar) / समाचार-पत्र}%
{Newspaper.}
\bhojexample{भोर में अखबार पढ़ल करीं।}{Read the newspaper in the morning.}

\bhojsubentry{अखबारवाला}{\K{\ba\bkha\bba\bmaa\bra\bva\bmaa\bla\bmaa}}{Akhbaarwala}{हि. पुं. | अरबी-तद्भव}{\starsThree}%
{अखबार बांटने वाला।}%
{अखबारवाला}%
{Newspaper vendor.}

\bhojentry{अखरा}{\K{\ba\bkha\bra\bmaa}}{Akhra}{हि. वि. | [ठेठ]}{\starsTwo}%
{बिना पॉलिश का, प्राकृतिक।}%
{अखरा (Akhra) / ठेठ}%
{Unpolished, raw, natural.}

\bhojentry{अखरोट}{\K{\ba\bkha\bra\bmo\btta}}{AkhroT}{हि. पुं. | तद्भव}{\starsTwo}%
{एक कड़ा पहाड़ी फल (सूखा मेवा)।}%
{अखरोट (AkhroT)}%
{Walnut fruit.}

\bhojentry{अखिरी}{\K{\ba\bkha\bmii\bra\bmii}}{Akhiri}{हि. वि. | अरबी आगत}{\starsThree}%
{अंत, अंतिम, समाप्ति।}%
{अखिरी (Akhiri) / अखीर}%
{Final, end, conclusion.}

\bhojentry{अगराइल}{\K{\ba\bga\bra\baa\bai\bla}}{Agril}{हि. वि. | [ठेठ]}{\starsTwo}%
{घमंडी, इतराया हुआ।}%
{अगराइल (Agril) / इतराइल}%
{Arrogant, haughty, pompous.}

\bhojentry{अगवा}{\K{\ba\bga\bva\bmaa}}{Agwa}{हि. पुं. | [ठेठ]}{\starsThree}%
{अगुआ, मार्गदर्शक, नेता।}%
{अगवा (Agwa) / नेता}%
{Leader, guide.}

\bhojentry{अगोरब}{\K{\ba\bga\bmo\bra\bba}}{Agorab}{हि. क्रि. | [ठेठ]}{\starsThree}%
{रखवाली करना, पहरा देना।}%
{अगोरब (Agorab) / पहरा}%
{To guard, watch over.}
\bhojexample{बाबा रात भर खेत अगोरले।}{Grandfather guarded the fields all night.}

\bhojsubentry{अगोरा}{\K{\ba\bga\bmo\bra\bmaa}}{Agora}{हि. पुं. | [ठेठ]}{\starsTwo}%
{रखवाली का काम या खेत की अगोरी।}%
{अगोरा (Agora) / पहरा}%
{Guarding, watchman duty.}

\bhojentry{अगहन}{\K{\ba\bga\bha\bna}}{Agahan}{हि. पुं. | तद्भव}{\starsThree}%
{नौवां हिंदू महीना (मार्गशीर्ष)।}%
{अगहन (Agahan)}%
{Margashirsha month.}

\bhojentry{अगली}{\K{\ba\bga\bla\bmii}}{Agali}{हि. वि. | तद्भव}{\starsThree}%
{सामने वाली, आने वाली।}%
{अगली (Agali) / अगिला}%
{Next, upcoming.}

\bhojentry{अगाड़ी}{\K{\ba\bga\bmaa\bdda\bnukta\bmii}}{Agaadi}{हि. स्त्री. | तद्भव}{\starsTwo}%
{सामने का भाग, अगला हिस्सा।}%
{अगाड़ी (Agaadi)}%
{Front part, forefront.}

\bhojentry{अगम}{\K{\ba\bga\bma}}{Agam}{सं. वि. | तत्सम}{\starsOne}%
{जहाँ पहुँचना कठिन हो, गहरा।}%
{अगम (Agam) / कठिन}%
{Inaccessible, deep.}

\bhojentry{अगुवाई}{\K{\ba\bga\bmu\bva\baa\bai}}{Aguwai}{हि. स्त्री. | तद्भव}{\starsThree}%
{नेतृत्व, राह दिखाना।}%
{अगुवाई (Aguwai) / आगू चलल}%
{Leadership, guidance.}

\bhojentry{अग्नि}{\K{\ba\bga\bvir\bna\bmi}}{Agni}{सं. स्त्री. | तत्सम}{\starsTwo}%
{आग, अनल, पावक।}%
{अग्नि (Agni) / आगि}%
{Fire.}
\bhojsee{आगि}

\bhojentry{अघाना}{\K{\ba\bgha\bmaa\bna\bmaa}}{Aghana}{हि. क्रि. | [ठेठ]}{\starsThree}%
{पेट भरना, तृप्त होना।}%
{अघाना (Aghana) / अघा जाइल}%
{To be satiated, full.}
\bhojexample{भोजन खा के लैका अघा गइल।}{The child was full after eating.}

\bhojsubentry{अघाएल}{\K{\ba\bgha\baa\bai\bla}}{Aghael}{हि. वि. | [ठेठ]}{\starsTwo}%
{तृप्त, जिसका पेट भरा हो।}%
{अघाएल}%
{Satiated.}

\bhojentry{अचरज}{\K{\ba\bca\bra\bja}}{Acaraj}{हि. पुं. | तद्भव}{\starsThree}%
{आश्चर्य, ताज्जुब।}%
{अचरज (Acaraj) / आश्चर्य}%
{Wonder, surprise, amazement.}

\bhojentry{अचार}{\K{\ba\bca\bmaa\bra}}{Acaar}{हि. पुं. | तद्भव}{\starsThree}%
{मसालेदार खट्टा खाद्य पदार्थ।}%
{अचार (Acaar)}%
{Pickle.}

\bhojentry{अचके}{\K{\ba\bca\bka\bme}}{Acake}{हि. अव्य. | [ठेठ]}{\starsThree}%
{एकाएक, अचानक।}%
{अचके (Acake) / अचानक}%
{Suddenly, abruptly.}
\bhojexample{अचके पानी आ गइल।}{Suddenly it started raining.}

\bhojentry{अचानक}{\K{\ba\bca\bmaa\bna\bka}}{Acaanak}{हि. अव्य. | फारसी आगत}{\starsThree}%
{बिना सूचना के, यकायक।}%
{अचानक (Acaanak) / अचके}%
{Suddenly, unexpected.}

\bhojentry{अचरा}{\K{\ba\bca\bra\bmaa}}{Acra}{हि. पुं. | [ठेठ]}{\starsThree}%
{साड़ी का पल्लू या आंचल।}%
{अचरा (Acra) / अँचरा}%
{Sari pallu, skirt border.}

\bhojentry{अचेत}{\K{\ba\bca\bme\bta}}{Acet}{हि. वि. | तद्भव}{\starsTwo}%
{बेहोश, सुध-बुध खोया हुआ।}%
{अचेत (Acet) / बेहोश}%
{Unconscious, faint.}

\bhojentry{अच्छा}{\K{\ba\bca\bvir\bcha\bmaa}}{Accha}{हि. वि. | तद्भव}{\starsThree}%
{सुंदर, नीक, बढ़िया।}%
{अच्छा (Accha) / नीक}%
{Good, fine, nice.}

\bhojentry{अटकल}{\K{\ba\btta\bka\bla}}{Atkal}{हि. स्त्री. | [ठेठ]}{\starsThree}%
{अनुमान, अंदाज।}%
{अटकल (Atkal) / अंटकल}%
{Conjecture, guess.}

\bhojentry{अटकब}{\K{\ba\btta\bka\bba}}{Atkab}{हि. क्रि. | [ठेठ]}{\starsThree}%
{फंसना, रुकना।}%
{अटकब (Atkab) / अटकल}%
{To get stuck, halted.}

\bhojentry{अटारी}{\K{\ba\btta\bmaa\bra\bmii}}{Ataari}{हि. स्त्री. | तद्भव}{\starsTwo}%
{कोठा, छत का कमरा।}%
{अटारी (Ataari) / कोठा}%
{Attic, balcony room.}

\bhojentry{अटिया}{\K{\ba\btta\bmi\bya\bmaa}}{Atiya}{हि. स्त्री. | [ठेठ]}{\starsTwo}%
{घास या सूत का छोटा बंडल।}%
{अटिया (Atiya) / अँकवारी}%
{Small sheaf, coil.}

\bhojentry{अटूट}{\K{\ba\btta\bmuu\btta}}{AtuuT}{हि. वि. | तद्भव}{\starsThree}%
{जो न टूटे, दृढ़।}%
{अटूट (AtuuT) / पक्का}%
{Unbreakable, firm.}

\bhojentry{अठारह}{\K{\ba\bttha\bmaa\bra\bha}}{AThaarah}{हि. सं. | तद्भव}{\starsThree}%
{१८ की संख्या।}%
{अठारह (AThaarah)}%
{Eighteen (18).}

\bhojentry{अठन्नी}{\K{\ba\bttha\bna\bvir\bna\bmii}}{AThannii}{हि. स्त्री. | तद्भव}{\starsTwo}%
{आठ आने का सिक्का।}%
{अठन्नी (AThannii)}%
{Eight-anna coin.}

\bhojentry{अड़ब}{\K{\ba\bdda\bnukta\bba}}{ARab}{हि. क्रि. | [ठेठ]}{\starsThree}%
{जिद्द करना, अपनी बात पर डटना।}%
{अड़ब (ARab) / डटल}%
{To insist, be stubborn.}

\bhojsubentry{अड़ंग}{\K{\ba\bdda\bnukta\banus\bga}}{ARang}{हि. पुं. | [ठेठ]}{\starsTwo}%
{बाधा, व्यवधान।}%
{अड़ंग / रुकावट}%
{Hindrance, obstacle.}

\bhojentry{अड़हुल}{\K{\ba\bdda\bnukta\bha\bmu\bla}}{ARhul}{हि. पुं. | [ठेठ]}{\starsThree}%
{गुड़हल का फूल।}%
{अड़हुल (ARhul) / गुड़हल}%
{Hibiscus flower.}

\bhojentry{अतनी}{\K{\ba\bta\bna\bmii}}{Atni}{हि. वि. | [ठेठ]}{\starsThree}%
{इतनी, बहुत छोटी मात्रा।}%
{अतनी (Atni) / तनी गो}%
{So small amount, this little.}

\bhojentry{अतवार}{\K{\ba\bta\bva\bmaa\bra}}{Atvaar}{हि. पुं. | तद्भव}{\starsThree}%
{रविवार का दिन।}%
{अतवार (Atvaar) / रविवार}%
{Sunday.}
\bhojexample{अतवार के स्कूल बंद रहेला।}{Schools remain closed on Sunday.}

\bhojentry{अतिथि}{\K{\ba\bta\bmi\btha\bmi}}{Atithi}{सं. पुं. | तत्सम}{\starsTwo}%
{पाहुन, मेहमान।}%
{पाहुन (Paahun) / अतिथि}%
{Guest, visitor.}

\bhojentry{अतर}{\K{\ba\bta\bra}}{Atar}{हि. पुं. | अरबी आगत}{\starsTwo}%
{इत्र, सुगंधि।}%
{इत्र (Itr) / अतर}%
{Perfume, attar.}

\bhojsubentry{अत्तरदानी}{\K{\ba\bta\bvir\bta\bra\bda\bmaa\bna\bmii}}{Attardaani}{हि. स्त्री. | अरबी-फारसी}{\starsOne}%
{इत्र रखने की डिब्बी।}%
{अत्तरदानी}%
{Perfume flask.}

\bhojentry{अथाह}{\K{\ba\btha\bmaa\bha}}{Athaah}{हि. वि. | तद्भव}{\starsTwo}%
{अगाध, जिसकी थाह न हो।}%
{अथाह (Athaah) / गहरा}%
{Fathomless, bottomless.}

\bhojentry{अदालत}{\K{\ba\bda\bmaa\bla\bta}}{Adaalat}{हि. स्त्री. | अरबी आगत}{\starsThree}%
{न्यायालय, कचेहरी।}%
{अदालत (Adaalat) / कचेहरी}%
{Court of law.}

\bhojentry{अदला-बदली}{\K{\ba\bda\bla\bmaa}}{Adla-badli}{हि. स्त्री. | [ठेठ]}{\starsThree}%
{अदला-बदली, विनिमय।}%
{अदला-बदली (Adla-badli)}%
{Barter, exchange.}

\bhojentry{अदवरी}{\K{\ba\bda\bva\bra\bmii}}{Advari}{हि. स्त्री. | [ठेठ]}{\starsThree}%
{दाल की सुखाई हुई बड़ी।}%
{अदवरी (Advari) / बड़ी}%
{Lentil dumplings (Badi).}

\bhojentry{अदरक}{\K{\ba\bda\bra\bka}}{Adrak}{हि. पुं. | तद्भव}{\starsThree}%
{आदी, एक कंद मसाला।}%
{आदी (Aadi) / अदरक}%
{Ginger.}
\bhojsee{आदी}

\bhojentry{अदहन}{\K{\ba\bda\bha\bna}}{Adahan}{हि. पुं. | [ठेठ]}{\starsThree}%
{चावल या दाल पकाने के लिए उबाला गया पानी।}%
{अदहन (Adahan)}%
{Boiled cooking water.}
\bhojexample{अदहन खौल गइल बा, चावल डाल दऽ।}{The water has boiled, add the rice.}

\bhojentry{अधिकार}{\K{\ba\bdha\bmi\bka\bmaa\bra}}{Adhikaar}{सं. पुं. | तत्सम}{\starsThree}%
{हक, आधिपत्य।}%
{अधिकार (Adhikaar) / हक}%
{Right, authority, power.}

\bhojsubentry{अधिकारी}{\K{\ba\bdha\bmi\bka\bmaa\bra\bmii}}{Adhikaarii}{सं. पुं. | तत्सम}{\starsThree}%
{अफसर, हकदार व्यक्ति।}%
{अधिकारी / अफसर}%
{Officer, official.}

\bhojentry{अनमोल}{\K{\ba\bna\bmo\bla}}{Anmol}{हि. वि. | तद्भव}{\starsThree}%
{बेकीमती, जिसका मूल्य न हो।}%
{अनमोल (Anmol) / बेकीमती}%
{Priceless, precious.}

\bhojentry{अनार}{\K{\ba\bna\bmaa\bra}}{Anaar}{हि. पुं. | फारसी आगत}{\starsThree}%
{एक प्रसिद्ध रसीला फल।}%
{अनार (Anaar)}%
{Pomegranate fruit.}

\bhojentry{अनाज}{\K{\ba\bna\bmaa\bja}}{Anaaj}{हि. पुं. | तद्भव}{\starsThree}%
{अन्न, खाद्यान्न।}%
{अनाज (Anaaj) / अन्न}%
{Food grain, cereal.}

\bhojentry{अनपढ़}{\K{\ba\bna\bpa\bddha\bnukta}}{AnpaDh}{हि. वि. | तद्भव}{\starsThree}%
{निरक्षर, बिना पढ़ा-लिखा।}%
{अनपढ़ (AnpaDh) / निरक्षर}%
{Illiterate, uneducated.}

\bhojentry{अनेक}{\K{\ba\bna\bme\bka}}{Anek}{सं. वि. | तत्सम}{\starsThree}%
{बहुत सारे, बहुविध।}%
{अनेक (Anek) / बहुतन}%
{Many, numerous.}

\bhojentry{अन्नदाता}{\K{\ba\bna\bvir\bna\bda\bmaa\bta\bmaa}}{Annadaata}{हि. पुं. | तद्भव}{\starsThree}%
{किसान, भोजन देने वाला।}%
{अन्नदाता (Annadaata) / किसान}%
{Provider of food, farmer.}

\bhojentry{अपराध}{\K{\ba\bpa\bra\bmaa\bdha}}{Aparaadh}{सं. पुं. | तत्सम}{\starsThree}%
{कसूर, जुर्म।}%
{अपराध (Aparaadh) / जुर्म}%
{Offense, crime.}

\bhojsubentry{अपराधी}{\K{\ba\bpa\bra\bmaa\bdha\bmii}}{Aparaadhii}{सं. पुं. | तत्सम}{\starsThree}%
{दोषी, मुजरिम।}%
{अपराधी / मुजरिम}%
{Culprit, offender.}

\bhojentry{अब}{\K{\ba\bba}}{Ab}{हि. अव्य. | तद्भव}{\starsThree}%
{इस समय, अबही।}%
{अब (Ab) / अबहीं}%
{Now, at present.}

\bhojsubentry{अबहीं}{\K{\ba\bba\bha\bmii\banus}}{Abhiiñ}{हि. अव्य. | [ठेठ]}{\starsThree}%
{इसी क्षण, तुरंत।}%
{अबहीं / तुरंत}%
{Right now, immediately.}

\bhojentry{अबेर}{\K{\ba\bba\bme\bra}}{Aber}{हि. स्त्री. | [ठेठ]}{\starsThree}%
{देरी, विलंब।}%
{अबेर (Aber) / देर}%
{Delay, late time.}

\bhojentry{अमरूद}{\K{\ba\bma\bra\bmuu\bda}}{Amruud}{हि. पुं. | फारसी आगत}{\starsThree}%
{एक मीठा फल (जाम/सूफड़ी)।}%
{अमरूद (Amruud) / सूफड़ी}%
{Guava fruit.}

\bhojentry{अमीर}{\K{\ba\bma\bmii\bra}}{Amiir}{हि. वि. | अरबी आगत}{\starsThree}%
{धनी, दौलतमंद।}%
{अमीर (Amiir) / धनी}%
{Rich, wealthy.}

\bhojentry{अरहर}{\K{\ba\bra\bha\bra}}{Arahar}{हि. स्त्री. | तद्भव/कृषि}{\starsThree}%
{तुवर दाल का पौधा या दाल।}%
{अरहर (Arahar) / राहर}%
{Pigeon pea (Tuvar dal).}
